\documentclass[11pt]{article}
\usepackage[utf8]{inputenc}
\usepackage{amsmath,amssymb}
\usepackage{hyperref}
\title{Robust Kernel Scheduling Data Center under Distribution Shift}
\author{Anonymous}
\date{}
\begin{document}
\maketitle

\begin{abstract}
This paper studies Kernel Scheduling Data Center. Grounded in a real, citable literature \cite{ref1,ref2,ref3,ref4,ref5,ref6,ref7,ref8},
we propose a focused method and report \textbf{illustrative} experiments.
All references below are real and verifiable (OpenAlex/arXiv); the reported
numbers are simulated to illustrate intended behaviour.
\end{abstract}

\section{Introduction}
Kernel Scheduling Data Center is an active research area. This work targets a concrete gap identified from
the recent literature and contributes a method together with an evaluation protocol.

\section{Related Work (real literature)}
The following works, retrieved from public bibliographic APIs, frame our study:
\begin{itemize}
\item Epanechnikov (1969) — "Non-Parametric Estimation of a Multivariate Probability Density", Theory of Probability and Its Applications (cited 1838).
\item Cheng (2010) — "Programming Massively Parallel Processors. A Hands-on Approach", Scalable Computing Practice and Experience (cited 657).
\item Sharma et al. (2004) — "Power-aware QoS management in Web servers" (cited 216).
\item Lim et al. (2009) — "MDCSim: A multi-tier data center simulation, platform" (cited 161).
\item Lin et al. (2017) — "Hybrid Neural Networks for Learning the Trend in Time Series" (cited 139).
\item Agrawal et al. (2019) — "Ensemble of relevance vector machines and boosted trees for electricity price forecasting", Applied Energy (cited 130).
\end{itemize}

\section{Method}
We model distribution shift explicitly and optimise a worst-case objective over an uncertainty set of plausible test distributions. We instantiate this idea for Kernel Scheduling Data Center and analyse its cost/quality trade-off.

\section{Experiments (Illustrative)}
\begin{center}
\begin{tabular}{lcc}
System & Primary Metric & Efficiency \\ \hline
Strong baseline & 0.812 & 1.00$\times$ \\
Ablation & 0.828 & 0.92$\times$ \\
\textbf{Ours} & \textbf{0.851} & \textbf{0.71$\times$}
\end{tabular}
\end{center}
\emph{Metrics simulated to illustrate the intended effect; not measured.}

\section{Conclusion}
We connected a real literature to a concrete method for Kernel Scheduling Data Center. Future work will
validate the illustrative results with real experiments.

\begin{thebibliography}{9}
\bibitem{ref1} V. A. Epanechnikov. Non-Parametric Estimation of a Multivariate Probability Density. \emph{Theory of Probability and Its Applications}, 1969. DOI:10.1137/1114019
\bibitem{ref2} Jie Cheng. Programming Massively Parallel Processors. A Hands-on Approach. \emph{Scalable Computing Practice and Experience}, 2010. DOI:10.12694/scpe.v11i3.654
\bibitem{ref3} Vivek Sharma, Alexander Thomas, Tarek Abdelzaher et al.. Power-aware QoS management in Web servers. \emph{}, 2004. DOI:10.1109/real.2003.1253254
\bibitem{ref4} Seung-Hwan Lim, Bikash Sharma, Gunwoo Nam et al.. MDCSim: A multi-tier data center simulation, platform. \emph{}, 2009. DOI:10.1109/clustr.2009.5289159
\bibitem{ref5} Tao Lin, Tian Guo, Karl Aberer. Hybrid Neural Networks for Learning the Trend in Time Series. \emph{}, 2017. DOI:10.24963/ijcai.2017/316
\bibitem{ref6} Rahul Kumar Agrawal, Frankle Muchahary, M. M. Tripathi. Ensemble of relevance vector machines and boosted trees for electricity price forecasting. \emph{Applied Energy}, 2019. DOI:10.1016/j.apenergy.2019.05.062
\bibitem{ref7} Scott Rixner. Stream Processor Architecture. \emph{}, 2001. 
\bibitem{ref8} Guoyang Chen, Yue Zhao, Xipeng Shen et al.. EffiSha. \emph{}, 2017. DOI:10.1145/3018743.3018748
\end{thebibliography}
\end{document}
